% Shared style for the ms2 paper figures (TikZ/pgfplots).
%
% Every fig-*.tex inputs this file, so a change here changes all figures.
% Data come from data/<fig>/ as written by
% bench/campaign-20261012/paper_figures_tikz.py.
%
% Load order: \documentclass, then % Shared style for the ms2 paper figures (TikZ/pgfplots).
%
% Every fig-*.tex inputs this file, so a change here changes all figures.
% Data come from data/<fig>/ as written by
% bench/campaign-20261012/paper_figures_tikz.py.
%
% Load order: \documentclass, then % Shared style for the ms2 paper figures (TikZ/pgfplots).
%
% Every fig-*.tex inputs this file, so a change here changes all figures.
% Data come from data/<fig>/ as written by
% bench/campaign-20261012/paper_figures_tikz.py.
%
% Load order: \documentclass, then % Shared style for the ms2 paper figures (TikZ/pgfplots).
%
% Every fig-*.tex inputs this file, so a change here changes all figures.
% Data come from data/<fig>/ as written by
% bench/campaign-20261012/paper_figures_tikz.py.
%
% Load order: \documentclass, then \input{figures-style}. The body font is
% newtxtext/newtxmath (Times-like), matching preview.tex.

\usepackage[T1]{fontenc}
\usepackage{newtxtext,newtxmath}
\usepackage{pgfplots}
\usepackage{pgfplotstable}
\usepgfplotslibrary{groupplots}
\usetikzlibrary{plotmarks,calc}
\pgfplotsset{compat=1.18}

% Colors. Layout colors, ropebwt3 red and the Others purple were validated as
% one set in paper_figures.py (no normal-vision failure). Psi aqua and dual
% dark yellow sit just under 3:1 contrast on white, so identity also rides on
% the legend and direct labels. Text never wears a series color.
\definecolor{figInk}{HTML}{0B0B0B}
\definecolor{figInkTwo}{HTML}{52514E}
\definecolor{figMuted}{HTML}{8A8984}
\definecolor{figGrid}{HTML}{E4E3DF}
\definecolor{msLcp}{HTML}{2A78D6}
\definecolor{msPsi}{HTML}{1BAF7A}
\definecolor{msDual}{HTML}{C98500}
\definecolor{ropeRed}{HTML}{E34948}
\definecolor{ropeLnineteen}{HTML}{4D4D4D}
\definecolor{ropeLtwentyfive}{HTML}{6E6E6E}
\definecolor{ropeLthirty}{HTML}{919191}
\definecolor{ropeLforty}{HTML}{B3B3B3}
\definecolor{othersPurple}{HTML}{7D3C98}
% ms2-phi appears only in Fig Ops. With the three layout colors it passes the
% palette checker; it is too close to othersPurple and ropeRed to share a
% figure with them.
\definecolor{msPhi}{HTML}{A3478F}

% Font sizes (absolute, so they do not depend on the class option). Figure
% text is sans serif (Helvetica clone from newtxtext's \sffamily) because
% option names such as -l1 read as -11 in Times.
\newcommand{\figTickFont}{\sffamily\fontsize{7}{8}\selectfont}
\newcommand{\figLabelFont}{\sffamily\fontsize{8}{9}\selectfont}
\newcommand{\figLegendFont}{\sffamily\fontsize{7.5}{8.5}\selectfont}
\newcommand{\figTitleFont}{\sffamily\fontsize{7.5}{8.5}\selectfont}
\newcommand{\figPointFont}{\sffamily\fontsize{6.5}{7.5}\selectfont}

% Panel geometry. A four-panel row: the left margin (tick labels plus y
% label) is about 0.32 in, so four panels and three gaps fill 6.5 in.
\newlength{\figPanelW}\setlength{\figPanelW}{1.38in}
\newlength{\figPanelH}\setlength{\figPanelH}{1.12in}
\newlength{\figPanelWthree}\setlength{\figPanelWthree}{1.88in}
\newlength{\figHsep}\setlength{\figHsep}{5pt}

% Marks that pgf's plotmarks library lacks: filled hexagon and filled star.
\pgfdeclareplotmark{hexagon*}{%
  \pgfpathmoveto{\pgfqpointpolar{90}{1.1\pgfplotmarksize}}%
  \foreach \a in {150,210,270,330,30}{\pgfpathlineto{\pgfqpointpolar{\a}{1.1\pgfplotmarksize}}}%
  \pgfpathclose\pgfusepathqfillstroke}
\pgfdeclareplotmark{fstar*}{%
  \pgfpathmoveto{\pgfqpointpolar{90}{1.45\pgfplotmarksize}}%
  \foreach \a/\r in {126/0.6,162/1.45,198/0.6,234/1.45,270/0.6,306/1.45,342/0.6,18/1.45,54/0.6}%
    {\pgfpathlineto{\pgfqpointpolar{\a}{\r\pgfplotmarksize}}}%
  \pgfpathclose\pgfusepathqfillstroke}
\pgfdeclareplotmark{fplus*}{%
  \pgfpathrectangle{\pgfqpoint{-1.25\pgfplotmarksize}{-0.4\pgfplotmarksize}}{\pgfqpoint{2.5\pgfplotmarksize}{0.8\pgfplotmarksize}}%
  \pgfpathrectangle{\pgfqpoint{-0.4\pgfplotmarksize}{-1.25\pgfplotmarksize}}{\pgfqpoint{0.8\pgfplotmarksize}{2.5\pgfplotmarksize}}%
  \pgfusepathqfill}

% Keeps only rows whose column #1 equals #2 (string compare).
\pgfplotsset{
  discard if not/.style 2 args={
    x filter/.append code={%
      \edef\figTempA{\thisrow{#1}}\edef\figTempB{#2}%
      \ifx\figTempA\figTempB\else\def\pgfmathresult{nan}\fi}},
  discard if/.style 2 args={
    x filter/.append code={%
      \edef\figTempA{\thisrow{#1}}\edef\figTempB{#2}%
      \ifx\figTempA\figTempB\def\pgfmathresult{nan}\fi}},
}

% Base axis style for every panel.
\pgfplotsset{
  fig axis/.style={
    scale only axis, width=\figPanelW, height=\figPanelH,
    axis lines*=left, minor tick length=0pt, axis line style={figMuted, line width=0.4pt},
    every tick/.style={figMuted, line width=0.4pt},
    major tick length=2pt,
    tick label style={font=\figTickFont, text=figInkTwo},
    label style={font=\figLabelFont, text=figInkTwo},
    ymajorgrids, grid style={figGrid, line width=0.35pt},
    xmin=0, xmax=\FigXmax, ymin=\FigYmin, ymax=\FigYmax,
    % ymode=log is given in each figure's axis options: inside a style it is
    % ignored by a plain axis environment (pgfplots reads it early).
    clip=false, unbounded coords=discard,
    xtick pos=bottom, ytick pos=left,
    xticklabel style={yshift=0.5pt}, yticklabel style={xshift=0.5pt},
    axis on top=false,
  },
  ms speed axis/.style={
    fig axis, ytick={0.01,0.1,1,10,100}, yticklabels={0.01,0.1,1,10,100},
    xtick={0,50,100,150,200}, xticklabels={0,,100,,200},
  },
  % Panels after the first drop the 0 label so it cannot touch the previous 200.
  ms speed rest/.style={xticklabels={,,100,,200}},
  chr22 speed axis/.style={
    fig axis, ytick={0.1,1,10,100}, yticklabels={0.1,1,10,100},
    xtick={0,0.5,1,1.5,2,2.5}, xticklabels={0,,1,,2,},
  },
  chr22 speed rest/.style={xticklabels={,,1,,2,}},
  locate axis/.style={
    fig axis, ytick={1e3,1e4,1e5,1e6,1e7}, yticklabels={1k,10k,100k,1M,10M},
    xtick={0,100,200,300}, xticklabels={0,100,200,300},
  },
  locate rest/.style={xticklabels={,100,200,300}},
  % Groupplot layout: tick labels only on the leftmost panel.
  % Used as group style={fig row=N}.
  group/fig row/.style={group size=#1 by 1, horizontal sep=\figHsep,
    yticklabels at=edge left, ylabels at=edge left},
  fig legend/.style={
    legend columns=-1,
    legend style={draw=none, fill=none, font=\figLegendFont, text=figInk,
      inner xsep=0pt, inner ysep=1pt,
      /tikz/every even column/.append style={column sep=5pt},
      /tikz/every odd column/.append style={column sep=1pt}},
    legend cell align=left,
  },
}
% Panel title: a small corner label inside the plot, top left.
\newcommand{\figPanelTitle}[1]{%
  \node[anchor=north west, inner sep=1.5pt, font=\figTitleFont, text=figInk]
    at (rel axis cs:0,1) {#1};}

% Series styles, one per data key. Filled marks carry a thin white rim so
% overlapping points stay separable.
\pgfplotsset{
  fig mark/.style={only marks, mark size=1.6pt, forget plot},
  fig filled/.style 2 args={fig mark, mark=#2,
    mark options={fill=#1, draw=white, line width=0.3pt, solid}},
  fig hollow/.style 2 args={fig mark, mark=#2,
    mark options={fill=white, draw=#1, line width=0.65pt, solid}},
  fig tick/.style 2 args={only marks, forget plot, mark=-, mark size=#2,
    mark options={draw=#1, line width=1.0pt, line cap=butt}},
  ms key/lcp/.style={fig filled={msLcp}{*}},
  ms key/psi/.style={fig filled={msPsi}{*}},
  ms key/dual/.style={fig filled={msDual}{*}},
  ms key/phi/.style={fig filled={msPhi}{*}},
  ms key/lcpF/.style={fig filled={msLcp}{diamond*}, mark size=2.3pt},
  ms key/psiF/.style={fig filled={msPsi}{diamond*}, mark size=2.3pt},
  ms key/phiF/.style={fig filled={msPhi}{diamond*}, mark size=2.3pt},
  ms key/dualF/.style={fig filled={msDual}{diamond*}, mark size=2.3pt},
  % Oracle replay: a hollow square in the layout's color, so it does not read
  % as an error-bar whisker.
  fig oracle/.style={fig hollow={#1}{square}, mark size=1.7pt, forget plot},
  ms key/lcpO/.style={fig oracle={msLcp}},
  ms key/psiO/.style={fig oracle={msPsi}},
  ms key/dualO/.style={fig oracle={msDual}},
  ms key/phiO/.style={fig oracle={msPhi}},
  ms key/rl1/.style={fig filled={ropeRed}{*}},
  ms key/rl19/.style={fig mark, mark=x, mark size=1.5pt,
    mark options={draw=ropeLnineteen, line width=0.8pt}},
  ms key/rl25/.style={fig mark, mark=triangle*,
    mark size=1.9pt, mark options={fill=ropeLtwentyfive, draw=white, line width=0.3pt, solid, rotate=180}},
  ms key/rl30/.style={fig filled={ropeLthirty}{pentagon*}},
  ms key/rl40/.style={fig filled={ropeLforty}{hexagon*}},
  ms key/moni/.style={fig hollow={othersPurple}{square}, mark size=1.5pt},
  ms key/phoni/.style={fig hollow={othersPurple}{diamond}, mark size=2.2pt},
  ms key/augphoni/.style={fig hollow={othersPurple}{triangle}, mark size=1.9pt},
  ms key/s11/.style={fig mark, mark=fplus*, mark size=1.7pt, mark options={fill=ropeRed}},
  ms key/s8/.style={fig filled={ropeRed}{fstar*}, mark size=2.0pt},
  ms key/s6/.style={fig filled={ropeRed}{hexagon*}, mark size=1.7pt},
  % Thin joining lines for rows that share a layout.
  fig link/.style={no marks, forget plot, line width=0.6pt, draw=#1},
  ms link/lcp/.style={fig link=msLcp},
  ms link/psi/.style={fig link=msPsi},
  ms link/dual/.style={fig link=msDual},
  ms link/phi/.style={fig link=msPhi},
  ms link/rs/.style={fig link=ropeRed!55},
  % Direct point labels from <panel>-labels.dat.
  fig point labels/.style={
    only marks, mark=none, forget plot,
    point meta=explicit symbolic,
    visualization depends on={value \thisrow{anchor}\as\figLabelAnchor},
    visualization depends on={\thisrow{dx}\as\figLabelDx},
    visualization depends on={\thisrow{dy}\as\figLabelDy},
    nodes near coords,
    nodes near coords style={font=\figPointFont, text=figInkTwo, inner sep=0.6pt,
      anchor=\figLabelAnchor, xshift=\figLabelDx pt, yshift=\figLabelDy pt},
  },
}

% \figPlotKeys{file}{key,key,...}: one \addplot per key, each marker shifted
% by the row's dx (pt). Keys absent from the file plot nothing.
\newcommand{\figPlotKeys}[2]{%
  \pgfplotsinvokeforeach{#2}{%
    \addplot[ms key/##1, discard if not={key}{##1},
      visualization depends on={\thisrow{dx}\as\figDx}, scatter,
      scatter/@pre marker code/.code={\pgfscope\pgftransformxshift{\figDx pt}},
      scatter/@post marker code/.code={\endpgfscope},
    ] table[x=x, y=y] {#1};}}
% \figPlotLinks{file}{layout,...}: joining lines, drawn first so marks sit on top.
\newcommand{\figPlotLinks}[2]{%
  \pgfplotsinvokeforeach{#2}{%
    \addplot[ms link/##1, discard if not={layout}{##1}] table[x=x, y=y] {#1};}}
\newcommand{\figPlotLabels}[1]{%
  \addplot[fig point labels] table[x=x, y=y, meta=text] {#1};}

% Short panel titles; the captions carry the full read-set descriptions.
\newcommand{\figTitle}[1]{\csname figtitle:#1\endcsname}
\expandafter\def\csname figtitle:human_short\endcsname{Short reads}
\expandafter\def\csname figtitle:human_pacbio\endcsname{PacBio HiFi}
\expandafter\def\csname figtitle:human_ont_modern\endcsname{Nanopore}
\expandafter\def\csname figtitle:guccione_skin\endcsname{Skin metagenome}
\expandafter\def\csname figtitle:simshort\endcsname{Sim.\ short reads}
\expandafter\def\csname figtitle:simpacbio\endcsname{Sim.\ PacBio HiFi}
\expandafter\def\csname figtitle:simont\endcsname{Sim.\ Nanopore}
\expandafter\def\csname figtitle:cap1\endcsname{\textit{N} = 1}
\expandafter\def\csname figtitle:cap10\endcsname{\textit{N} = 10}
\expandafter\def\csname figtitle:cap100\endcsname{\textit{N} = 100}

% \figPanelBody{fig}{panel}{link layouts}{keys}: everything drawn inside one
% panel, after its \nextgroupplot.
\newcommand{\figPanelBody}[4]{%
  \figPlotLinks{data/#1/#2.dat}{#3}%
  \figPlotKeys{data/#1/#2.dat}{#4}%
  \figPlotLabels{data/#1/#2-labels.dat}%
  \figPanelTitle{\figTitle{#2}}}

% Legend groups. An empty-legend entry names the tool before its modes.
\newcommand{\figLegendGroupName}[1]{\addlegendimage{empty legend}\addlegendentry{\textit{#1}}}
\newcommand{\figLegendKey}[2]{\addlegendimage{ms key/#1}\addlegendentry{#2}}
\newcommand{\figLegendMS}{%
  \figLegendGroupName{ms2:}%
  \figLegendKey{lcp}{min}\figLegendKey{psi}{psi}\figLegendKey{dual}{dual}%
  \figLegendGroupName{ropebwt3:}%
  \figLegendKey{rl1}{-l1}\figLegendKey{rl19}{-l19*}\figLegendKey{rl25}{-l25*}%
  \figLegendKey{rl30}{-l30*}\figLegendKey{rl40}{-l40*}}
\newcommand{\figLegendOracle}{%
  \addlegendimage{fig hollow={figInkTwo}{square}, mark size=1.7pt}\addlegendentry{oracle}}
\newcommand{\figLegendFtab}{%
  \addlegendimage{fig filled={figInkTwo}{*}}\addlegendentry{ordinary}%
  \addlegendimage{fig filled={figInkTwo}{diamond*}, mark size=2.3pt}\addlegendentry{ftab}}
\newcommand{\figLegendOthers}{%
  \figLegendKey{moni}{MONI}\figLegendKey{phoni}{PHONI}\figLegendKey{augphoni}{AugPHONI}}
\newcommand{\figLegendLocate}{%
  \figLegendGroupName{ms2:}%
  \figLegendKey{psi}{psi}\figLegendKey{dual}{dual}%
  \figLegendGroupName{ropebwt3 -l1:}%
  \figLegendKey{s11}{-s11}\figLegendKey{s8}{-s8}\figLegendKey{s6}{-s6}}

% Shared x label and legend placed relative to a row of panels.
% \figRowFurniture{first}{last}{x label}{legend name}
\newcommand{\figRowFurniture}[4]{%
  \node[font=\figLabelFont, text=figInkTwo, anchor=north]
    at ($(#1.south west)!0.5!(#2.south east)+(0,-8pt)$) {#3};
  \node[anchor=south, inner sep=0pt]
    at ($(#1.north west)!0.5!(#2.north east)+(0,2pt)$) {\pgfplotslegendfromname{#4}};}

\newcommand{\figMemLabel}{Peak resident memory (GiB)}
\newcommand{\figSpeedLabel}{Speed (Mbases/s)}
\newcommand{\figLocLabel}{Locations/s}

% Fig Ops: bar panels for time and operations per base (data/fig-ops/).
% Bars are drawn per layout key with zero bar shift so each sits on its own x.
% An ops bar is the full total in the layout color with the LF segment (one
% step per base, which is also what the oracle performs) painted over it as a
% 45 percent tint; the white rim leaves a surface gap between the segments.
\newlength{\figOpsPanelW}\setlength{\figOpsPanelW}{1.25in}
\newlength{\figOpsCountPanelW}\setlength{\figOpsCountPanelW}{1.0in}
\pgfplotsset{
  ops axis/.style={
    fig axis, width=\figOpsPanelW, xmin=-0.6, xmax=3.6, ymin=0,
    xtick={0,1,2,3}, xticklabels={min,psi,phi,dual}, major x tick style={draw=none},
    xticklabel style={font=\figLabelFont, text=figInkTwo},
    % Numeric tick labels outside math mode, so they take the sans tick font.
    yticklabel style={/pgf/number format/assume math mode=true},
  },
  % Operation panels leave out ms2-min (keyed lcp), whose only added work is a short LCP scan.
  ops count axis/.style={ops axis, width=\figOpsCountPanelW, xmin=0.4, xmax=3.6,
    xtick={1,2,3}, xticklabels={psi,phi,dual}},
  ops bar/.style={ybar, bar shift=0pt, bar width=0.62, forget plot, draw=white, line width=0.5pt, fill=#1},
  ops tick/.style={only marks, forget plot, mark=-, mark size=8.6pt,
    mark options={draw=figInk, line width=1.0pt, line cap=butt}},
  ops labels/.style={only marks, mark=none, forget plot, point meta=explicit symbolic,
    nodes near coords, nodes near coords style={font=\figPointFont, text=figInkTwo, inner sep=0.6pt,
      anchor=south, yshift=0.6pt}},
  ops legend bar/.style={legend image code/.code={\fill[#1] (0pt,-2.6pt) rectangle (11pt,2.6pt);}},
}
\newcommand{\figOpsColor}[1]{\csname figopscolor:#1\endcsname}
\expandafter\def\csname figopscolor:lcp\endcsname{msLcp}
\expandafter\def\csname figopscolor:psi\endcsname{msPsi}
\expandafter\def\csname figopscolor:dual\endcsname{msDual}
\expandafter\def\csname figopscolor:phi\endcsname{msPhi}
% \figOpsTime{file}: ordinary time bars, oracle ticks and bar labels.
\newcommand{\figOpsTime}[1]{%
  \pgfplotsinvokeforeach{lcp,psi,dual,phi}{%
    \addplot[ops bar=\figOpsColor{##1}, discard if not={key}{##1}] table[x=x, y=ns] {#1};}%
  \addplot[ops tick] table[x=x, y=oracle] {#1};
  \addplot[ops labels] table[x=x, y=ns, meta=nslab] {#1};}
% \figOpsOps{file}: stacked operation bars and total labels.
\newcommand{\figOpsOps}[1]{%
  \pgfplotsinvokeforeach{psi,dual,phi}{%
    \addplot[ops bar=\figOpsColor{##1}, discard if not={key}{##1}] table[x=x, y=total] {#1};
    \addplot[ops bar=\figOpsColor{##1}!45, discard if not={key}{##1}] table[x=x, y=lf] {#1};}%
  \addplot[ops labels, discard if={key}{lcp}] table[x=x, y=total, meta=totlab] {#1};}
\newcommand{\figLegendOps}{%
  \addlegendimage{ops legend bar=msLcp}\addlegendentry{min}%
  \addlegendimage{ops legend bar=msPsi}\addlegendentry{psi}%
  \addlegendimage{ops legend bar=msPhi}\addlegendentry{phi}%
  \addlegendimage{ops legend bar=msDual}\addlegendentry{dual}%
  \addlegendimage{line legend, draw=figInk, line width=1.0pt}\addlegendentry{oracle time}%
  \addlegendimage{ops legend bar=figInkTwo!45}\addlegendentry{LF steps}%
  \addlegendimage{ops legend bar=figInkTwo}\addlegendentry{Recovery ($\psi$ \& $\phi$) steps}}
% Bar panels put the title above the plot, clear of the bar labels.
\newcommand{\figOpsTitle}[1]{%
  \node[anchor=south west, inner sep=1.5pt, font=\figTitleFont, text=figInk]
    at (rel axis cs:0,1) {#1};}
\newcommand{\figOpsTimeLabel}{ns per base (query time)}
\newcommand{\figOpsLabel}{Operations per base}
\expandafter\def\csname figtitle:mtb152\endcsname{mtb152}
\expandafter\def\csname figtitle:human472\endcsname{HPRC2}
. The body font is
% newtxtext/newtxmath (Times-like), matching preview.tex.

\usepackage[T1]{fontenc}
\usepackage{newtxtext,newtxmath}
\usepackage{pgfplots}
\usepackage{pgfplotstable}
\usepgfplotslibrary{groupplots}
\usetikzlibrary{plotmarks,calc}
\pgfplotsset{compat=1.18}

% Colors. Layout colors, ropebwt3 red and the Others purple were validated as
% one set in paper_figures.py (no normal-vision failure). Psi aqua and dual
% dark yellow sit just under 3:1 contrast on white, so identity also rides on
% the legend and direct labels. Text never wears a series color.
\definecolor{figInk}{HTML}{0B0B0B}
\definecolor{figInkTwo}{HTML}{52514E}
\definecolor{figMuted}{HTML}{8A8984}
\definecolor{figGrid}{HTML}{E4E3DF}
\definecolor{msLcp}{HTML}{2A78D6}
\definecolor{msPsi}{HTML}{1BAF7A}
\definecolor{msDual}{HTML}{C98500}
\definecolor{ropeRed}{HTML}{E34948}
\definecolor{ropeLnineteen}{HTML}{4D4D4D}
\definecolor{ropeLtwentyfive}{HTML}{6E6E6E}
\definecolor{ropeLthirty}{HTML}{919191}
\definecolor{ropeLforty}{HTML}{B3B3B3}
\definecolor{othersPurple}{HTML}{7D3C98}
% ms2-phi appears only in Fig Ops. With the three layout colors it passes the
% palette checker; it is too close to othersPurple and ropeRed to share a
% figure with them.
\definecolor{msPhi}{HTML}{A3478F}

% Font sizes (absolute, so they do not depend on the class option). Figure
% text is sans serif (Helvetica clone from newtxtext's \sffamily) because
% option names such as -l1 read as -11 in Times.
\newcommand{\figTickFont}{\sffamily\fontsize{7}{8}\selectfont}
\newcommand{\figLabelFont}{\sffamily\fontsize{8}{9}\selectfont}
\newcommand{\figLegendFont}{\sffamily\fontsize{7.5}{8.5}\selectfont}
\newcommand{\figTitleFont}{\sffamily\fontsize{7.5}{8.5}\selectfont}
\newcommand{\figPointFont}{\sffamily\fontsize{6.5}{7.5}\selectfont}

% Panel geometry. A four-panel row: the left margin (tick labels plus y
% label) is about 0.32 in, so four panels and three gaps fill 6.5 in.
\newlength{\figPanelW}\setlength{\figPanelW}{1.38in}
\newlength{\figPanelH}\setlength{\figPanelH}{1.12in}
\newlength{\figPanelWthree}\setlength{\figPanelWthree}{1.88in}
\newlength{\figHsep}\setlength{\figHsep}{5pt}

% Marks that pgf's plotmarks library lacks: filled hexagon and filled star.
\pgfdeclareplotmark{hexagon*}{%
  \pgfpathmoveto{\pgfqpointpolar{90}{1.1\pgfplotmarksize}}%
  \foreach \a in {150,210,270,330,30}{\pgfpathlineto{\pgfqpointpolar{\a}{1.1\pgfplotmarksize}}}%
  \pgfpathclose\pgfusepathqfillstroke}
\pgfdeclareplotmark{fstar*}{%
  \pgfpathmoveto{\pgfqpointpolar{90}{1.45\pgfplotmarksize}}%
  \foreach \a/\r in {126/0.6,162/1.45,198/0.6,234/1.45,270/0.6,306/1.45,342/0.6,18/1.45,54/0.6}%
    {\pgfpathlineto{\pgfqpointpolar{\a}{\r\pgfplotmarksize}}}%
  \pgfpathclose\pgfusepathqfillstroke}
\pgfdeclareplotmark{fplus*}{%
  \pgfpathrectangle{\pgfqpoint{-1.25\pgfplotmarksize}{-0.4\pgfplotmarksize}}{\pgfqpoint{2.5\pgfplotmarksize}{0.8\pgfplotmarksize}}%
  \pgfpathrectangle{\pgfqpoint{-0.4\pgfplotmarksize}{-1.25\pgfplotmarksize}}{\pgfqpoint{0.8\pgfplotmarksize}{2.5\pgfplotmarksize}}%
  \pgfusepathqfill}

% Keeps only rows whose column #1 equals #2 (string compare).
\pgfplotsset{
  discard if not/.style 2 args={
    x filter/.append code={%
      \edef\figTempA{\thisrow{#1}}\edef\figTempB{#2}%
      \ifx\figTempA\figTempB\else\def\pgfmathresult{nan}\fi}},
  discard if/.style 2 args={
    x filter/.append code={%
      \edef\figTempA{\thisrow{#1}}\edef\figTempB{#2}%
      \ifx\figTempA\figTempB\def\pgfmathresult{nan}\fi}},
}

% Base axis style for every panel.
\pgfplotsset{
  fig axis/.style={
    scale only axis, width=\figPanelW, height=\figPanelH,
    axis lines*=left, minor tick length=0pt, axis line style={figMuted, line width=0.4pt},
    every tick/.style={figMuted, line width=0.4pt},
    major tick length=2pt,
    tick label style={font=\figTickFont, text=figInkTwo},
    label style={font=\figLabelFont, text=figInkTwo},
    ymajorgrids, grid style={figGrid, line width=0.35pt},
    xmin=0, xmax=\FigXmax, ymin=\FigYmin, ymax=\FigYmax,
    % ymode=log is given in each figure's axis options: inside a style it is
    % ignored by a plain axis environment (pgfplots reads it early).
    clip=false, unbounded coords=discard,
    xtick pos=bottom, ytick pos=left,
    xticklabel style={yshift=0.5pt}, yticklabel style={xshift=0.5pt},
    axis on top=false,
  },
  ms speed axis/.style={
    fig axis, ytick={0.01,0.1,1,10,100}, yticklabels={0.01,0.1,1,10,100},
    xtick={0,50,100,150,200}, xticklabels={0,,100,,200},
  },
  % Panels after the first drop the 0 label so it cannot touch the previous 200.
  ms speed rest/.style={xticklabels={,,100,,200}},
  chr22 speed axis/.style={
    fig axis, ytick={0.1,1,10,100}, yticklabels={0.1,1,10,100},
    xtick={0,0.5,1,1.5,2,2.5}, xticklabels={0,,1,,2,},
  },
  chr22 speed rest/.style={xticklabels={,,1,,2,}},
  locate axis/.style={
    fig axis, ytick={1e3,1e4,1e5,1e6,1e7}, yticklabels={1k,10k,100k,1M,10M},
    xtick={0,100,200,300}, xticklabels={0,100,200,300},
  },
  locate rest/.style={xticklabels={,100,200,300}},
  % Groupplot layout: tick labels only on the leftmost panel.
  % Used as group style={fig row=N}.
  group/fig row/.style={group size=#1 by 1, horizontal sep=\figHsep,
    yticklabels at=edge left, ylabels at=edge left},
  fig legend/.style={
    legend columns=-1,
    legend style={draw=none, fill=none, font=\figLegendFont, text=figInk,
      inner xsep=0pt, inner ysep=1pt,
      /tikz/every even column/.append style={column sep=5pt},
      /tikz/every odd column/.append style={column sep=1pt}},
    legend cell align=left,
  },
}
% Panel title: a small corner label inside the plot, top left.
\newcommand{\figPanelTitle}[1]{%
  \node[anchor=north west, inner sep=1.5pt, font=\figTitleFont, text=figInk]
    at (rel axis cs:0,1) {#1};}

% Series styles, one per data key. Filled marks carry a thin white rim so
% overlapping points stay separable.
\pgfplotsset{
  fig mark/.style={only marks, mark size=1.6pt, forget plot},
  fig filled/.style 2 args={fig mark, mark=#2,
    mark options={fill=#1, draw=white, line width=0.3pt, solid}},
  fig hollow/.style 2 args={fig mark, mark=#2,
    mark options={fill=white, draw=#1, line width=0.65pt, solid}},
  fig tick/.style 2 args={only marks, forget plot, mark=-, mark size=#2,
    mark options={draw=#1, line width=1.0pt, line cap=butt}},
  ms key/lcp/.style={fig filled={msLcp}{*}},
  ms key/psi/.style={fig filled={msPsi}{*}},
  ms key/dual/.style={fig filled={msDual}{*}},
  ms key/phi/.style={fig filled={msPhi}{*}},
  ms key/lcpF/.style={fig filled={msLcp}{diamond*}, mark size=2.3pt},
  ms key/psiF/.style={fig filled={msPsi}{diamond*}, mark size=2.3pt},
  ms key/phiF/.style={fig filled={msPhi}{diamond*}, mark size=2.3pt},
  ms key/dualF/.style={fig filled={msDual}{diamond*}, mark size=2.3pt},
  % Oracle replay: a hollow square in the layout's color, so it does not read
  % as an error-bar whisker.
  fig oracle/.style={fig hollow={#1}{square}, mark size=1.7pt, forget plot},
  ms key/lcpO/.style={fig oracle={msLcp}},
  ms key/psiO/.style={fig oracle={msPsi}},
  ms key/dualO/.style={fig oracle={msDual}},
  ms key/phiO/.style={fig oracle={msPhi}},
  ms key/rl1/.style={fig filled={ropeRed}{*}},
  ms key/rl19/.style={fig mark, mark=x, mark size=1.5pt,
    mark options={draw=ropeLnineteen, line width=0.8pt}},
  ms key/rl25/.style={fig mark, mark=triangle*,
    mark size=1.9pt, mark options={fill=ropeLtwentyfive, draw=white, line width=0.3pt, solid, rotate=180}},
  ms key/rl30/.style={fig filled={ropeLthirty}{pentagon*}},
  ms key/rl40/.style={fig filled={ropeLforty}{hexagon*}},
  ms key/moni/.style={fig hollow={othersPurple}{square}, mark size=1.5pt},
  ms key/phoni/.style={fig hollow={othersPurple}{diamond}, mark size=2.2pt},
  ms key/augphoni/.style={fig hollow={othersPurple}{triangle}, mark size=1.9pt},
  ms key/s11/.style={fig mark, mark=fplus*, mark size=1.7pt, mark options={fill=ropeRed}},
  ms key/s8/.style={fig filled={ropeRed}{fstar*}, mark size=2.0pt},
  ms key/s6/.style={fig filled={ropeRed}{hexagon*}, mark size=1.7pt},
  % Thin joining lines for rows that share a layout.
  fig link/.style={no marks, forget plot, line width=0.6pt, draw=#1},
  ms link/lcp/.style={fig link=msLcp},
  ms link/psi/.style={fig link=msPsi},
  ms link/dual/.style={fig link=msDual},
  ms link/phi/.style={fig link=msPhi},
  ms link/rs/.style={fig link=ropeRed!55},
  % Direct point labels from <panel>-labels.dat.
  fig point labels/.style={
    only marks, mark=none, forget plot,
    point meta=explicit symbolic,
    visualization depends on={value \thisrow{anchor}\as\figLabelAnchor},
    visualization depends on={\thisrow{dx}\as\figLabelDx},
    visualization depends on={\thisrow{dy}\as\figLabelDy},
    nodes near coords,
    nodes near coords style={font=\figPointFont, text=figInkTwo, inner sep=0.6pt,
      anchor=\figLabelAnchor, xshift=\figLabelDx pt, yshift=\figLabelDy pt},
  },
}

% \figPlotKeys{file}{key,key,...}: one \addplot per key, each marker shifted
% by the row's dx (pt). Keys absent from the file plot nothing.
\newcommand{\figPlotKeys}[2]{%
  \pgfplotsinvokeforeach{#2}{%
    \addplot[ms key/##1, discard if not={key}{##1},
      visualization depends on={\thisrow{dx}\as\figDx}, scatter,
      scatter/@pre marker code/.code={\pgfscope\pgftransformxshift{\figDx pt}},
      scatter/@post marker code/.code={\endpgfscope},
    ] table[x=x, y=y] {#1};}}
% \figPlotLinks{file}{layout,...}: joining lines, drawn first so marks sit on top.
\newcommand{\figPlotLinks}[2]{%
  \pgfplotsinvokeforeach{#2}{%
    \addplot[ms link/##1, discard if not={layout}{##1}] table[x=x, y=y] {#1};}}
\newcommand{\figPlotLabels}[1]{%
  \addplot[fig point labels] table[x=x, y=y, meta=text] {#1};}

% Short panel titles; the captions carry the full read-set descriptions.
\newcommand{\figTitle}[1]{\csname figtitle:#1\endcsname}
\expandafter\def\csname figtitle:human_short\endcsname{Short reads}
\expandafter\def\csname figtitle:human_pacbio\endcsname{PacBio HiFi}
\expandafter\def\csname figtitle:human_ont_modern\endcsname{Nanopore}
\expandafter\def\csname figtitle:guccione_skin\endcsname{Skin metagenome}
\expandafter\def\csname figtitle:simshort\endcsname{Sim.\ short reads}
\expandafter\def\csname figtitle:simpacbio\endcsname{Sim.\ PacBio HiFi}
\expandafter\def\csname figtitle:simont\endcsname{Sim.\ Nanopore}
\expandafter\def\csname figtitle:cap1\endcsname{\textit{N} = 1}
\expandafter\def\csname figtitle:cap10\endcsname{\textit{N} = 10}
\expandafter\def\csname figtitle:cap100\endcsname{\textit{N} = 100}

% \figPanelBody{fig}{panel}{link layouts}{keys}: everything drawn inside one
% panel, after its \nextgroupplot.
\newcommand{\figPanelBody}[4]{%
  \figPlotLinks{data/#1/#2.dat}{#3}%
  \figPlotKeys{data/#1/#2.dat}{#4}%
  \figPlotLabels{data/#1/#2-labels.dat}%
  \figPanelTitle{\figTitle{#2}}}

% Legend groups. An empty-legend entry names the tool before its modes.
\newcommand{\figLegendGroupName}[1]{\addlegendimage{empty legend}\addlegendentry{\textit{#1}}}
\newcommand{\figLegendKey}[2]{\addlegendimage{ms key/#1}\addlegendentry{#2}}
\newcommand{\figLegendMS}{%
  \figLegendGroupName{ms2:}%
  \figLegendKey{lcp}{min}\figLegendKey{psi}{psi}\figLegendKey{dual}{dual}%
  \figLegendGroupName{ropebwt3:}%
  \figLegendKey{rl1}{-l1}\figLegendKey{rl19}{-l19*}\figLegendKey{rl25}{-l25*}%
  \figLegendKey{rl30}{-l30*}\figLegendKey{rl40}{-l40*}}
\newcommand{\figLegendOracle}{%
  \addlegendimage{fig hollow={figInkTwo}{square}, mark size=1.7pt}\addlegendentry{oracle}}
\newcommand{\figLegendFtab}{%
  \addlegendimage{fig filled={figInkTwo}{*}}\addlegendentry{ordinary}%
  \addlegendimage{fig filled={figInkTwo}{diamond*}, mark size=2.3pt}\addlegendentry{ftab}}
\newcommand{\figLegendOthers}{%
  \figLegendKey{moni}{MONI}\figLegendKey{phoni}{PHONI}\figLegendKey{augphoni}{AugPHONI}}
\newcommand{\figLegendLocate}{%
  \figLegendGroupName{ms2:}%
  \figLegendKey{psi}{psi}\figLegendKey{dual}{dual}%
  \figLegendGroupName{ropebwt3 -l1:}%
  \figLegendKey{s11}{-s11}\figLegendKey{s8}{-s8}\figLegendKey{s6}{-s6}}

% Shared x label and legend placed relative to a row of panels.
% \figRowFurniture{first}{last}{x label}{legend name}
\newcommand{\figRowFurniture}[4]{%
  \node[font=\figLabelFont, text=figInkTwo, anchor=north]
    at ($(#1.south west)!0.5!(#2.south east)+(0,-8pt)$) {#3};
  \node[anchor=south, inner sep=0pt]
    at ($(#1.north west)!0.5!(#2.north east)+(0,2pt)$) {\pgfplotslegendfromname{#4}};}

\newcommand{\figMemLabel}{Peak resident memory (GiB)}
\newcommand{\figSpeedLabel}{Speed (Mbases/s)}
\newcommand{\figLocLabel}{Locations/s}

% Fig Ops: bar panels for time and operations per base (data/fig-ops/).
% Bars are drawn per layout key with zero bar shift so each sits on its own x.
% An ops bar is the full total in the layout color with the LF segment (one
% step per base, which is also what the oracle performs) painted over it as a
% 45 percent tint; the white rim leaves a surface gap between the segments.
\newlength{\figOpsPanelW}\setlength{\figOpsPanelW}{1.25in}
\newlength{\figOpsCountPanelW}\setlength{\figOpsCountPanelW}{1.0in}
\pgfplotsset{
  ops axis/.style={
    fig axis, width=\figOpsPanelW, xmin=-0.6, xmax=3.6, ymin=0,
    xtick={0,1,2,3}, xticklabels={min,psi,phi,dual}, major x tick style={draw=none},
    xticklabel style={font=\figLabelFont, text=figInkTwo},
    % Numeric tick labels outside math mode, so they take the sans tick font.
    yticklabel style={/pgf/number format/assume math mode=true},
  },
  % Operation panels leave out ms2-min (keyed lcp), whose only added work is a short LCP scan.
  ops count axis/.style={ops axis, width=\figOpsCountPanelW, xmin=0.4, xmax=3.6,
    xtick={1,2,3}, xticklabels={psi,phi,dual}},
  ops bar/.style={ybar, bar shift=0pt, bar width=0.62, forget plot, draw=white, line width=0.5pt, fill=#1},
  ops tick/.style={only marks, forget plot, mark=-, mark size=8.6pt,
    mark options={draw=figInk, line width=1.0pt, line cap=butt}},
  ops labels/.style={only marks, mark=none, forget plot, point meta=explicit symbolic,
    nodes near coords, nodes near coords style={font=\figPointFont, text=figInkTwo, inner sep=0.6pt,
      anchor=south, yshift=0.6pt}},
  ops legend bar/.style={legend image code/.code={\fill[#1] (0pt,-2.6pt) rectangle (11pt,2.6pt);}},
}
\newcommand{\figOpsColor}[1]{\csname figopscolor:#1\endcsname}
\expandafter\def\csname figopscolor:lcp\endcsname{msLcp}
\expandafter\def\csname figopscolor:psi\endcsname{msPsi}
\expandafter\def\csname figopscolor:dual\endcsname{msDual}
\expandafter\def\csname figopscolor:phi\endcsname{msPhi}
% \figOpsTime{file}: ordinary time bars, oracle ticks and bar labels.
\newcommand{\figOpsTime}[1]{%
  \pgfplotsinvokeforeach{lcp,psi,dual,phi}{%
    \addplot[ops bar=\figOpsColor{##1}, discard if not={key}{##1}] table[x=x, y=ns] {#1};}%
  \addplot[ops tick] table[x=x, y=oracle] {#1};
  \addplot[ops labels] table[x=x, y=ns, meta=nslab] {#1};}
% \figOpsOps{file}: stacked operation bars and total labels.
\newcommand{\figOpsOps}[1]{%
  \pgfplotsinvokeforeach{psi,dual,phi}{%
    \addplot[ops bar=\figOpsColor{##1}, discard if not={key}{##1}] table[x=x, y=total] {#1};
    \addplot[ops bar=\figOpsColor{##1}!45, discard if not={key}{##1}] table[x=x, y=lf] {#1};}%
  \addplot[ops labels, discard if={key}{lcp}] table[x=x, y=total, meta=totlab] {#1};}
\newcommand{\figLegendOps}{%
  \addlegendimage{ops legend bar=msLcp}\addlegendentry{min}%
  \addlegendimage{ops legend bar=msPsi}\addlegendentry{psi}%
  \addlegendimage{ops legend bar=msPhi}\addlegendentry{phi}%
  \addlegendimage{ops legend bar=msDual}\addlegendentry{dual}%
  \addlegendimage{line legend, draw=figInk, line width=1.0pt}\addlegendentry{oracle time}%
  \addlegendimage{ops legend bar=figInkTwo!45}\addlegendentry{LF steps}%
  \addlegendimage{ops legend bar=figInkTwo}\addlegendentry{Recovery ($\psi$ \& $\phi$) steps}}
% Bar panels put the title above the plot, clear of the bar labels.
\newcommand{\figOpsTitle}[1]{%
  \node[anchor=south west, inner sep=1.5pt, font=\figTitleFont, text=figInk]
    at (rel axis cs:0,1) {#1};}
\newcommand{\figOpsTimeLabel}{ns per base (query time)}
\newcommand{\figOpsLabel}{Operations per base}
\expandafter\def\csname figtitle:mtb152\endcsname{mtb152}
\expandafter\def\csname figtitle:human472\endcsname{HPRC2}
. The body font is
% newtxtext/newtxmath (Times-like), matching preview.tex.

\usepackage[T1]{fontenc}
\usepackage{newtxtext,newtxmath}
\usepackage{pgfplots}
\usepackage{pgfplotstable}
\usepgfplotslibrary{groupplots}
\usetikzlibrary{plotmarks,calc}
\pgfplotsset{compat=1.18}

% Colors. Layout colors, ropebwt3 red and the Others purple were validated as
% one set in paper_figures.py (no normal-vision failure). Psi aqua and dual
% dark yellow sit just under 3:1 contrast on white, so identity also rides on
% the legend and direct labels. Text never wears a series color.
\definecolor{figInk}{HTML}{0B0B0B}
\definecolor{figInkTwo}{HTML}{52514E}
\definecolor{figMuted}{HTML}{8A8984}
\definecolor{figGrid}{HTML}{E4E3DF}
\definecolor{msLcp}{HTML}{2A78D6}
\definecolor{msPsi}{HTML}{1BAF7A}
\definecolor{msDual}{HTML}{C98500}
\definecolor{ropeRed}{HTML}{E34948}
\definecolor{ropeLnineteen}{HTML}{4D4D4D}
\definecolor{ropeLtwentyfive}{HTML}{6E6E6E}
\definecolor{ropeLthirty}{HTML}{919191}
\definecolor{ropeLforty}{HTML}{B3B3B3}
\definecolor{othersPurple}{HTML}{7D3C98}
% ms2-phi appears only in Fig Ops. With the three layout colors it passes the
% palette checker; it is too close to othersPurple and ropeRed to share a
% figure with them.
\definecolor{msPhi}{HTML}{A3478F}

% Font sizes (absolute, so they do not depend on the class option). Figure
% text is sans serif (Helvetica clone from newtxtext's \sffamily) because
% option names such as -l1 read as -11 in Times.
\newcommand{\figTickFont}{\sffamily\fontsize{7}{8}\selectfont}
\newcommand{\figLabelFont}{\sffamily\fontsize{8}{9}\selectfont}
\newcommand{\figLegendFont}{\sffamily\fontsize{7.5}{8.5}\selectfont}
\newcommand{\figTitleFont}{\sffamily\fontsize{7.5}{8.5}\selectfont}
\newcommand{\figPointFont}{\sffamily\fontsize{6.5}{7.5}\selectfont}

% Panel geometry. A four-panel row: the left margin (tick labels plus y
% label) is about 0.32 in, so four panels and three gaps fill 6.5 in.
\newlength{\figPanelW}\setlength{\figPanelW}{1.38in}
\newlength{\figPanelH}\setlength{\figPanelH}{1.12in}
\newlength{\figPanelWthree}\setlength{\figPanelWthree}{1.88in}
\newlength{\figHsep}\setlength{\figHsep}{5pt}

% Marks that pgf's plotmarks library lacks: filled hexagon and filled star.
\pgfdeclareplotmark{hexagon*}{%
  \pgfpathmoveto{\pgfqpointpolar{90}{1.1\pgfplotmarksize}}%
  \foreach \a in {150,210,270,330,30}{\pgfpathlineto{\pgfqpointpolar{\a}{1.1\pgfplotmarksize}}}%
  \pgfpathclose\pgfusepathqfillstroke}
\pgfdeclareplotmark{fstar*}{%
  \pgfpathmoveto{\pgfqpointpolar{90}{1.45\pgfplotmarksize}}%
  \foreach \a/\r in {126/0.6,162/1.45,198/0.6,234/1.45,270/0.6,306/1.45,342/0.6,18/1.45,54/0.6}%
    {\pgfpathlineto{\pgfqpointpolar{\a}{\r\pgfplotmarksize}}}%
  \pgfpathclose\pgfusepathqfillstroke}
\pgfdeclareplotmark{fplus*}{%
  \pgfpathrectangle{\pgfqpoint{-1.25\pgfplotmarksize}{-0.4\pgfplotmarksize}}{\pgfqpoint{2.5\pgfplotmarksize}{0.8\pgfplotmarksize}}%
  \pgfpathrectangle{\pgfqpoint{-0.4\pgfplotmarksize}{-1.25\pgfplotmarksize}}{\pgfqpoint{0.8\pgfplotmarksize}{2.5\pgfplotmarksize}}%
  \pgfusepathqfill}

% Keeps only rows whose column #1 equals #2 (string compare).
\pgfplotsset{
  discard if not/.style 2 args={
    x filter/.append code={%
      \edef\figTempA{\thisrow{#1}}\edef\figTempB{#2}%
      \ifx\figTempA\figTempB\else\def\pgfmathresult{nan}\fi}},
  discard if/.style 2 args={
    x filter/.append code={%
      \edef\figTempA{\thisrow{#1}}\edef\figTempB{#2}%
      \ifx\figTempA\figTempB\def\pgfmathresult{nan}\fi}},
}

% Base axis style for every panel.
\pgfplotsset{
  fig axis/.style={
    scale only axis, width=\figPanelW, height=\figPanelH,
    axis lines*=left, minor tick length=0pt, axis line style={figMuted, line width=0.4pt},
    every tick/.style={figMuted, line width=0.4pt},
    major tick length=2pt,
    tick label style={font=\figTickFont, text=figInkTwo},
    label style={font=\figLabelFont, text=figInkTwo},
    ymajorgrids, grid style={figGrid, line width=0.35pt},
    xmin=0, xmax=\FigXmax, ymin=\FigYmin, ymax=\FigYmax,
    % ymode=log is given in each figure's axis options: inside a style it is
    % ignored by a plain axis environment (pgfplots reads it early).
    clip=false, unbounded coords=discard,
    xtick pos=bottom, ytick pos=left,
    xticklabel style={yshift=0.5pt}, yticklabel style={xshift=0.5pt},
    axis on top=false,
  },
  ms speed axis/.style={
    fig axis, ytick={0.01,0.1,1,10,100}, yticklabels={0.01,0.1,1,10,100},
    xtick={0,50,100,150,200}, xticklabels={0,,100,,200},
  },
  % Panels after the first drop the 0 label so it cannot touch the previous 200.
  ms speed rest/.style={xticklabels={,,100,,200}},
  chr22 speed axis/.style={
    fig axis, ytick={0.1,1,10,100}, yticklabels={0.1,1,10,100},
    xtick={0,0.5,1,1.5,2,2.5}, xticklabels={0,,1,,2,},
  },
  chr22 speed rest/.style={xticklabels={,,1,,2,}},
  locate axis/.style={
    fig axis, ytick={1e3,1e4,1e5,1e6,1e7}, yticklabels={1k,10k,100k,1M,10M},
    xtick={0,100,200,300}, xticklabels={0,100,200,300},
  },
  locate rest/.style={xticklabels={,100,200,300}},
  % Groupplot layout: tick labels only on the leftmost panel.
  % Used as group style={fig row=N}.
  group/fig row/.style={group size=#1 by 1, horizontal sep=\figHsep,
    yticklabels at=edge left, ylabels at=edge left},
  fig legend/.style={
    legend columns=-1,
    legend style={draw=none, fill=none, font=\figLegendFont, text=figInk,
      inner xsep=0pt, inner ysep=1pt,
      /tikz/every even column/.append style={column sep=5pt},
      /tikz/every odd column/.append style={column sep=1pt}},
    legend cell align=left,
  },
}
% Panel title: a small corner label inside the plot, top left.
\newcommand{\figPanelTitle}[1]{%
  \node[anchor=north west, inner sep=1.5pt, font=\figTitleFont, text=figInk]
    at (rel axis cs:0,1) {#1};}

% Series styles, one per data key. Filled marks carry a thin white rim so
% overlapping points stay separable.
\pgfplotsset{
  fig mark/.style={only marks, mark size=1.6pt, forget plot},
  fig filled/.style 2 args={fig mark, mark=#2,
    mark options={fill=#1, draw=white, line width=0.3pt, solid}},
  fig hollow/.style 2 args={fig mark, mark=#2,
    mark options={fill=white, draw=#1, line width=0.65pt, solid}},
  fig tick/.style 2 args={only marks, forget plot, mark=-, mark size=#2,
    mark options={draw=#1, line width=1.0pt, line cap=butt}},
  ms key/lcp/.style={fig filled={msLcp}{*}},
  ms key/psi/.style={fig filled={msPsi}{*}},
  ms key/dual/.style={fig filled={msDual}{*}},
  ms key/phi/.style={fig filled={msPhi}{*}},
  ms key/lcpF/.style={fig filled={msLcp}{diamond*}, mark size=2.3pt},
  ms key/psiF/.style={fig filled={msPsi}{diamond*}, mark size=2.3pt},
  ms key/phiF/.style={fig filled={msPhi}{diamond*}, mark size=2.3pt},
  ms key/dualF/.style={fig filled={msDual}{diamond*}, mark size=2.3pt},
  % Oracle replay: a hollow square in the layout's color, so it does not read
  % as an error-bar whisker.
  fig oracle/.style={fig hollow={#1}{square}, mark size=1.7pt, forget plot},
  ms key/lcpO/.style={fig oracle={msLcp}},
  ms key/psiO/.style={fig oracle={msPsi}},
  ms key/dualO/.style={fig oracle={msDual}},
  ms key/phiO/.style={fig oracle={msPhi}},
  ms key/rl1/.style={fig filled={ropeRed}{*}},
  ms key/rl19/.style={fig mark, mark=x, mark size=1.5pt,
    mark options={draw=ropeLnineteen, line width=0.8pt}},
  ms key/rl25/.style={fig mark, mark=triangle*,
    mark size=1.9pt, mark options={fill=ropeLtwentyfive, draw=white, line width=0.3pt, solid, rotate=180}},
  ms key/rl30/.style={fig filled={ropeLthirty}{pentagon*}},
  ms key/rl40/.style={fig filled={ropeLforty}{hexagon*}},
  ms key/moni/.style={fig hollow={othersPurple}{square}, mark size=1.5pt},
  ms key/phoni/.style={fig hollow={othersPurple}{diamond}, mark size=2.2pt},
  ms key/augphoni/.style={fig hollow={othersPurple}{triangle}, mark size=1.9pt},
  ms key/s11/.style={fig mark, mark=fplus*, mark size=1.7pt, mark options={fill=ropeRed}},
  ms key/s8/.style={fig filled={ropeRed}{fstar*}, mark size=2.0pt},
  ms key/s6/.style={fig filled={ropeRed}{hexagon*}, mark size=1.7pt},
  % Thin joining lines for rows that share a layout.
  fig link/.style={no marks, forget plot, line width=0.6pt, draw=#1},
  ms link/lcp/.style={fig link=msLcp},
  ms link/psi/.style={fig link=msPsi},
  ms link/dual/.style={fig link=msDual},
  ms link/phi/.style={fig link=msPhi},
  ms link/rs/.style={fig link=ropeRed!55},
  % Direct point labels from <panel>-labels.dat.
  fig point labels/.style={
    only marks, mark=none, forget plot,
    point meta=explicit symbolic,
    visualization depends on={value \thisrow{anchor}\as\figLabelAnchor},
    visualization depends on={\thisrow{dx}\as\figLabelDx},
    visualization depends on={\thisrow{dy}\as\figLabelDy},
    nodes near coords,
    nodes near coords style={font=\figPointFont, text=figInkTwo, inner sep=0.6pt,
      anchor=\figLabelAnchor, xshift=\figLabelDx pt, yshift=\figLabelDy pt},
  },
}

% \figPlotKeys{file}{key,key,...}: one \addplot per key, each marker shifted
% by the row's dx (pt). Keys absent from the file plot nothing.
\newcommand{\figPlotKeys}[2]{%
  \pgfplotsinvokeforeach{#2}{%
    \addplot[ms key/##1, discard if not={key}{##1},
      visualization depends on={\thisrow{dx}\as\figDx}, scatter,
      scatter/@pre marker code/.code={\pgfscope\pgftransformxshift{\figDx pt}},
      scatter/@post marker code/.code={\endpgfscope},
    ] table[x=x, y=y] {#1};}}
% \figPlotLinks{file}{layout,...}: joining lines, drawn first so marks sit on top.
\newcommand{\figPlotLinks}[2]{%
  \pgfplotsinvokeforeach{#2}{%
    \addplot[ms link/##1, discard if not={layout}{##1}] table[x=x, y=y] {#1};}}
\newcommand{\figPlotLabels}[1]{%
  \addplot[fig point labels] table[x=x, y=y, meta=text] {#1};}

% Short panel titles; the captions carry the full read-set descriptions.
\newcommand{\figTitle}[1]{\csname figtitle:#1\endcsname}
\expandafter\def\csname figtitle:human_short\endcsname{Short reads}
\expandafter\def\csname figtitle:human_pacbio\endcsname{PacBio HiFi}
\expandafter\def\csname figtitle:human_ont_modern\endcsname{Nanopore}
\expandafter\def\csname figtitle:guccione_skin\endcsname{Skin metagenome}
\expandafter\def\csname figtitle:simshort\endcsname{Sim.\ short reads}
\expandafter\def\csname figtitle:simpacbio\endcsname{Sim.\ PacBio HiFi}
\expandafter\def\csname figtitle:simont\endcsname{Sim.\ Nanopore}
\expandafter\def\csname figtitle:cap1\endcsname{\textit{N} = 1}
\expandafter\def\csname figtitle:cap10\endcsname{\textit{N} = 10}
\expandafter\def\csname figtitle:cap100\endcsname{\textit{N} = 100}

% \figPanelBody{fig}{panel}{link layouts}{keys}: everything drawn inside one
% panel, after its \nextgroupplot.
\newcommand{\figPanelBody}[4]{%
  \figPlotLinks{data/#1/#2.dat}{#3}%
  \figPlotKeys{data/#1/#2.dat}{#4}%
  \figPlotLabels{data/#1/#2-labels.dat}%
  \figPanelTitle{\figTitle{#2}}}

% Legend groups. An empty-legend entry names the tool before its modes.
\newcommand{\figLegendGroupName}[1]{\addlegendimage{empty legend}\addlegendentry{\textit{#1}}}
\newcommand{\figLegendKey}[2]{\addlegendimage{ms key/#1}\addlegendentry{#2}}
\newcommand{\figLegendMS}{%
  \figLegendGroupName{ms2:}%
  \figLegendKey{lcp}{min}\figLegendKey{psi}{psi}\figLegendKey{dual}{dual}%
  \figLegendGroupName{ropebwt3:}%
  \figLegendKey{rl1}{-l1}\figLegendKey{rl19}{-l19*}\figLegendKey{rl25}{-l25*}%
  \figLegendKey{rl30}{-l30*}\figLegendKey{rl40}{-l40*}}
\newcommand{\figLegendOracle}{%
  \addlegendimage{fig hollow={figInkTwo}{square}, mark size=1.7pt}\addlegendentry{oracle}}
\newcommand{\figLegendFtab}{%
  \addlegendimage{fig filled={figInkTwo}{*}}\addlegendentry{ordinary}%
  \addlegendimage{fig filled={figInkTwo}{diamond*}, mark size=2.3pt}\addlegendentry{ftab}}
\newcommand{\figLegendOthers}{%
  \figLegendKey{moni}{MONI}\figLegendKey{phoni}{PHONI}\figLegendKey{augphoni}{AugPHONI}}
\newcommand{\figLegendLocate}{%
  \figLegendGroupName{ms2:}%
  \figLegendKey{psi}{psi}\figLegendKey{dual}{dual}%
  \figLegendGroupName{ropebwt3 -l1:}%
  \figLegendKey{s11}{-s11}\figLegendKey{s8}{-s8}\figLegendKey{s6}{-s6}}

% Shared x label and legend placed relative to a row of panels.
% \figRowFurniture{first}{last}{x label}{legend name}
\newcommand{\figRowFurniture}[4]{%
  \node[font=\figLabelFont, text=figInkTwo, anchor=north]
    at ($(#1.south west)!0.5!(#2.south east)+(0,-8pt)$) {#3};
  \node[anchor=south, inner sep=0pt]
    at ($(#1.north west)!0.5!(#2.north east)+(0,2pt)$) {\pgfplotslegendfromname{#4}};}

\newcommand{\figMemLabel}{Peak resident memory (GiB)}
\newcommand{\figSpeedLabel}{Speed (Mbases/s)}
\newcommand{\figLocLabel}{Locations/s}

% Fig Ops: bar panels for time and operations per base (data/fig-ops/).
% Bars are drawn per layout key with zero bar shift so each sits on its own x.
% An ops bar is the full total in the layout color with the LF segment (one
% step per base, which is also what the oracle performs) painted over it as a
% 45 percent tint; the white rim leaves a surface gap between the segments.
\newlength{\figOpsPanelW}\setlength{\figOpsPanelW}{1.25in}
\newlength{\figOpsCountPanelW}\setlength{\figOpsCountPanelW}{1.0in}
\pgfplotsset{
  ops axis/.style={
    fig axis, width=\figOpsPanelW, xmin=-0.6, xmax=3.6, ymin=0,
    xtick={0,1,2,3}, xticklabels={min,psi,phi,dual}, major x tick style={draw=none},
    xticklabel style={font=\figLabelFont, text=figInkTwo},
    % Numeric tick labels outside math mode, so they take the sans tick font.
    yticklabel style={/pgf/number format/assume math mode=true},
  },
  % Operation panels leave out ms2-min (keyed lcp), whose only added work is a short LCP scan.
  ops count axis/.style={ops axis, width=\figOpsCountPanelW, xmin=0.4, xmax=3.6,
    xtick={1,2,3}, xticklabels={psi,phi,dual}},
  ops bar/.style={ybar, bar shift=0pt, bar width=0.62, forget plot, draw=white, line width=0.5pt, fill=#1},
  ops tick/.style={only marks, forget plot, mark=-, mark size=8.6pt,
    mark options={draw=figInk, line width=1.0pt, line cap=butt}},
  ops labels/.style={only marks, mark=none, forget plot, point meta=explicit symbolic,
    nodes near coords, nodes near coords style={font=\figPointFont, text=figInkTwo, inner sep=0.6pt,
      anchor=south, yshift=0.6pt}},
  ops legend bar/.style={legend image code/.code={\fill[#1] (0pt,-2.6pt) rectangle (11pt,2.6pt);}},
}
\newcommand{\figOpsColor}[1]{\csname figopscolor:#1\endcsname}
\expandafter\def\csname figopscolor:lcp\endcsname{msLcp}
\expandafter\def\csname figopscolor:psi\endcsname{msPsi}
\expandafter\def\csname figopscolor:dual\endcsname{msDual}
\expandafter\def\csname figopscolor:phi\endcsname{msPhi}
% \figOpsTime{file}: ordinary time bars, oracle ticks and bar labels.
\newcommand{\figOpsTime}[1]{%
  \pgfplotsinvokeforeach{lcp,psi,dual,phi}{%
    \addplot[ops bar=\figOpsColor{##1}, discard if not={key}{##1}] table[x=x, y=ns] {#1};}%
  \addplot[ops tick] table[x=x, y=oracle] {#1};
  \addplot[ops labels] table[x=x, y=ns, meta=nslab] {#1};}
% \figOpsOps{file}: stacked operation bars and total labels.
\newcommand{\figOpsOps}[1]{%
  \pgfplotsinvokeforeach{psi,dual,phi}{%
    \addplot[ops bar=\figOpsColor{##1}, discard if not={key}{##1}] table[x=x, y=total] {#1};
    \addplot[ops bar=\figOpsColor{##1}!45, discard if not={key}{##1}] table[x=x, y=lf] {#1};}%
  \addplot[ops labels, discard if={key}{lcp}] table[x=x, y=total, meta=totlab] {#1};}
\newcommand{\figLegendOps}{%
  \addlegendimage{ops legend bar=msLcp}\addlegendentry{min}%
  \addlegendimage{ops legend bar=msPsi}\addlegendentry{psi}%
  \addlegendimage{ops legend bar=msPhi}\addlegendentry{phi}%
  \addlegendimage{ops legend bar=msDual}\addlegendentry{dual}%
  \addlegendimage{line legend, draw=figInk, line width=1.0pt}\addlegendentry{oracle time}%
  \addlegendimage{ops legend bar=figInkTwo!45}\addlegendentry{LF steps}%
  \addlegendimage{ops legend bar=figInkTwo}\addlegendentry{Recovery ($\psi$ \& $\phi$) steps}}
% Bar panels put the title above the plot, clear of the bar labels.
\newcommand{\figOpsTitle}[1]{%
  \node[anchor=south west, inner sep=1.5pt, font=\figTitleFont, text=figInk]
    at (rel axis cs:0,1) {#1};}
\newcommand{\figOpsTimeLabel}{ns per base (query time)}
\newcommand{\figOpsLabel}{Operations per base}
\expandafter\def\csname figtitle:mtb152\endcsname{mtb152}
\expandafter\def\csname figtitle:human472\endcsname{HPRC2}
. The body font is
% newtxtext/newtxmath (Times-like), matching preview.tex.

\usepackage[T1]{fontenc}
\usepackage{newtxtext,newtxmath}
\usepackage{pgfplots}
\usepackage{pgfplotstable}
\usepgfplotslibrary{groupplots}
\usetikzlibrary{plotmarks,calc}
\pgfplotsset{compat=1.18}

% Colors. Layout colors, ropebwt3 red and the Others purple were validated as
% one set in paper_figures.py (no normal-vision failure). Psi aqua and dual
% dark yellow sit just under 3:1 contrast on white, so identity also rides on
% the legend and direct labels. Text never wears a series color.
\definecolor{figInk}{HTML}{0B0B0B}
\definecolor{figInkTwo}{HTML}{52514E}
\definecolor{figMuted}{HTML}{8A8984}
\definecolor{figGrid}{HTML}{E4E3DF}
\definecolor{msLcp}{HTML}{2A78D6}
\definecolor{msPsi}{HTML}{1BAF7A}
\definecolor{msDual}{HTML}{C98500}
\definecolor{ropeRed}{HTML}{E34948}
\definecolor{ropeLnineteen}{HTML}{4D4D4D}
\definecolor{ropeLtwentyfive}{HTML}{6E6E6E}
\definecolor{ropeLthirty}{HTML}{919191}
\definecolor{ropeLforty}{HTML}{B3B3B3}
\definecolor{othersPurple}{HTML}{7D3C98}
% ms2-phi appears only in Fig Ops. With the three layout colors it passes the
% palette checker; it is too close to othersPurple and ropeRed to share a
% figure with them.
\definecolor{msPhi}{HTML}{A3478F}

% Font sizes (absolute, so they do not depend on the class option). Figure
% text is sans serif (Helvetica clone from newtxtext's \sffamily) because
% option names such as -l1 read as -11 in Times.
\newcommand{\figTickFont}{\sffamily\fontsize{7}{8}\selectfont}
\newcommand{\figLabelFont}{\sffamily\fontsize{8}{9}\selectfont}
\newcommand{\figLegendFont}{\sffamily\fontsize{7.5}{8.5}\selectfont}
\newcommand{\figTitleFont}{\sffamily\fontsize{7.5}{8.5}\selectfont}
\newcommand{\figPointFont}{\sffamily\fontsize{6.5}{7.5}\selectfont}

% Panel geometry. A four-panel row: the left margin (tick labels plus y
% label) is about 0.32 in, so four panels and three gaps fill 6.5 in.
\newlength{\figPanelW}\setlength{\figPanelW}{1.38in}
\newlength{\figPanelH}\setlength{\figPanelH}{1.12in}
\newlength{\figPanelWthree}\setlength{\figPanelWthree}{1.88in}
\newlength{\figHsep}\setlength{\figHsep}{5pt}

% Marks that pgf's plotmarks library lacks: filled hexagon and filled star.
\pgfdeclareplotmark{hexagon*}{%
  \pgfpathmoveto{\pgfqpointpolar{90}{1.1\pgfplotmarksize}}%
  \foreach \a in {150,210,270,330,30}{\pgfpathlineto{\pgfqpointpolar{\a}{1.1\pgfplotmarksize}}}%
  \pgfpathclose\pgfusepathqfillstroke}
\pgfdeclareplotmark{fstar*}{%
  \pgfpathmoveto{\pgfqpointpolar{90}{1.45\pgfplotmarksize}}%
  \foreach \a/\r in {126/0.6,162/1.45,198/0.6,234/1.45,270/0.6,306/1.45,342/0.6,18/1.45,54/0.6}%
    {\pgfpathlineto{\pgfqpointpolar{\a}{\r\pgfplotmarksize}}}%
  \pgfpathclose\pgfusepathqfillstroke}
\pgfdeclareplotmark{fplus*}{%
  \pgfpathrectangle{\pgfqpoint{-1.25\pgfplotmarksize}{-0.4\pgfplotmarksize}}{\pgfqpoint{2.5\pgfplotmarksize}{0.8\pgfplotmarksize}}%
  \pgfpathrectangle{\pgfqpoint{-0.4\pgfplotmarksize}{-1.25\pgfplotmarksize}}{\pgfqpoint{0.8\pgfplotmarksize}{2.5\pgfplotmarksize}}%
  \pgfusepathqfill}

% Keeps only rows whose column #1 equals #2 (string compare).
\pgfplotsset{
  discard if not/.style 2 args={
    x filter/.append code={%
      \edef\figTempA{\thisrow{#1}}\edef\figTempB{#2}%
      \ifx\figTempA\figTempB\else\def\pgfmathresult{nan}\fi}},
  discard if/.style 2 args={
    x filter/.append code={%
      \edef\figTempA{\thisrow{#1}}\edef\figTempB{#2}%
      \ifx\figTempA\figTempB\def\pgfmathresult{nan}\fi}},
}

% Base axis style for every panel.
\pgfplotsset{
  fig axis/.style={
    scale only axis, width=\figPanelW, height=\figPanelH,
    axis lines*=left, minor tick length=0pt, axis line style={figMuted, line width=0.4pt},
    every tick/.style={figMuted, line width=0.4pt},
    major tick length=2pt,
    tick label style={font=\figTickFont, text=figInkTwo},
    label style={font=\figLabelFont, text=figInkTwo},
    ymajorgrids, grid style={figGrid, line width=0.35pt},
    xmin=0, xmax=\FigXmax, ymin=\FigYmin, ymax=\FigYmax,
    % ymode=log is given in each figure's axis options: inside a style it is
    % ignored by a plain axis environment (pgfplots reads it early).
    clip=false, unbounded coords=discard,
    xtick pos=bottom, ytick pos=left,
    xticklabel style={yshift=0.5pt}, yticklabel style={xshift=0.5pt},
    axis on top=false,
  },
  ms speed axis/.style={
    fig axis, ytick={0.01,0.1,1,10,100}, yticklabels={0.01,0.1,1,10,100},
    xtick={0,50,100,150,200}, xticklabels={0,,100,,200},
  },
  % Panels after the first drop the 0 label so it cannot touch the previous 200.
  ms speed rest/.style={xticklabels={,,100,,200}},
  chr22 speed axis/.style={
    fig axis, ytick={0.1,1,10,100}, yticklabels={0.1,1,10,100},
    xtick={0,0.5,1,1.5,2,2.5}, xticklabels={0,,1,,2,},
  },
  chr22 speed rest/.style={xticklabels={,,1,,2,}},
  locate axis/.style={
    fig axis, ytick={1e3,1e4,1e5,1e6,1e7}, yticklabels={1k,10k,100k,1M,10M},
    xtick={0,100,200,300}, xticklabels={0,100,200,300},
  },
  locate rest/.style={xticklabels={,100,200,300}},
  % Groupplot layout: tick labels only on the leftmost panel.
  % Used as group style={fig row=N}.
  group/fig row/.style={group size=#1 by 1, horizontal sep=\figHsep,
    yticklabels at=edge left, ylabels at=edge left},
  fig legend/.style={
    legend columns=-1,
    legend style={draw=none, fill=none, font=\figLegendFont, text=figInk,
      inner xsep=0pt, inner ysep=1pt,
      /tikz/every even column/.append style={column sep=5pt},
      /tikz/every odd column/.append style={column sep=1pt}},
    legend cell align=left,
  },
}
% Panel title: a small corner label inside the plot, top left.
\newcommand{\figPanelTitle}[1]{%
  \node[anchor=north west, inner sep=1.5pt, font=\figTitleFont, text=figInk]
    at (rel axis cs:0,1) {#1};}

% Series styles, one per data key. Filled marks carry a thin white rim so
% overlapping points stay separable.
\pgfplotsset{
  fig mark/.style={only marks, mark size=1.6pt, forget plot},
  fig filled/.style 2 args={fig mark, mark=#2,
    mark options={fill=#1, draw=white, line width=0.3pt, solid}},
  fig hollow/.style 2 args={fig mark, mark=#2,
    mark options={fill=white, draw=#1, line width=0.65pt, solid}},
  fig tick/.style 2 args={only marks, forget plot, mark=-, mark size=#2,
    mark options={draw=#1, line width=1.0pt, line cap=butt}},
  ms key/lcp/.style={fig filled={msLcp}{*}},
  ms key/psi/.style={fig filled={msPsi}{*}},
  ms key/dual/.style={fig filled={msDual}{*}},
  ms key/phi/.style={fig filled={msPhi}{*}},
  ms key/lcpF/.style={fig filled={msLcp}{diamond*}, mark size=2.3pt},
  ms key/psiF/.style={fig filled={msPsi}{diamond*}, mark size=2.3pt},
  ms key/phiF/.style={fig filled={msPhi}{diamond*}, mark size=2.3pt},
  ms key/dualF/.style={fig filled={msDual}{diamond*}, mark size=2.3pt},
  % Oracle replay: a hollow square in the layout's color, so it does not read
  % as an error-bar whisker.
  fig oracle/.style={fig hollow={#1}{square}, mark size=1.7pt, forget plot},
  ms key/lcpO/.style={fig oracle={msLcp}},
  ms key/psiO/.style={fig oracle={msPsi}},
  ms key/dualO/.style={fig oracle={msDual}},
  ms key/phiO/.style={fig oracle={msPhi}},
  ms key/rl1/.style={fig filled={ropeRed}{*}},
  ms key/rl19/.style={fig mark, mark=x, mark size=1.5pt,
    mark options={draw=ropeLnineteen, line width=0.8pt}},
  ms key/rl25/.style={fig mark, mark=triangle*,
    mark size=1.9pt, mark options={fill=ropeLtwentyfive, draw=white, line width=0.3pt, solid, rotate=180}},
  ms key/rl30/.style={fig filled={ropeLthirty}{pentagon*}},
  ms key/rl40/.style={fig filled={ropeLforty}{hexagon*}},
  ms key/moni/.style={fig hollow={othersPurple}{square}, mark size=1.5pt},
  ms key/phoni/.style={fig hollow={othersPurple}{diamond}, mark size=2.2pt},
  ms key/augphoni/.style={fig hollow={othersPurple}{triangle}, mark size=1.9pt},
  ms key/s11/.style={fig mark, mark=fplus*, mark size=1.7pt, mark options={fill=ropeRed}},
  ms key/s8/.style={fig filled={ropeRed}{fstar*}, mark size=2.0pt},
  ms key/s6/.style={fig filled={ropeRed}{hexagon*}, mark size=1.7pt},
  % Thin joining lines for rows that share a layout.
  fig link/.style={no marks, forget plot, line width=0.6pt, draw=#1},
  ms link/lcp/.style={fig link=msLcp},
  ms link/psi/.style={fig link=msPsi},
  ms link/dual/.style={fig link=msDual},
  ms link/phi/.style={fig link=msPhi},
  ms link/rs/.style={fig link=ropeRed!55},
  % Direct point labels from <panel>-labels.dat.
  fig point labels/.style={
    only marks, mark=none, forget plot,
    point meta=explicit symbolic,
    visualization depends on={value \thisrow{anchor}\as\figLabelAnchor},
    visualization depends on={\thisrow{dx}\as\figLabelDx},
    visualization depends on={\thisrow{dy}\as\figLabelDy},
    nodes near coords,
    nodes near coords style={font=\figPointFont, text=figInkTwo, inner sep=0.6pt,
      anchor=\figLabelAnchor, xshift=\figLabelDx pt, yshift=\figLabelDy pt},
  },
}

% \figPlotKeys{file}{key,key,...}: one \addplot per key, each marker shifted
% by the row's dx (pt). Keys absent from the file plot nothing.
\newcommand{\figPlotKeys}[2]{%
  \pgfplotsinvokeforeach{#2}{%
    \addplot[ms key/##1, discard if not={key}{##1},
      visualization depends on={\thisrow{dx}\as\figDx}, scatter,
      scatter/@pre marker code/.code={\pgfscope\pgftransformxshift{\figDx pt}},
      scatter/@post marker code/.code={\endpgfscope},
    ] table[x=x, y=y] {#1};}}
% \figPlotLinks{file}{layout,...}: joining lines, drawn first so marks sit on top.
\newcommand{\figPlotLinks}[2]{%
  \pgfplotsinvokeforeach{#2}{%
    \addplot[ms link/##1, discard if not={layout}{##1}] table[x=x, y=y] {#1};}}
\newcommand{\figPlotLabels}[1]{%
  \addplot[fig point labels] table[x=x, y=y, meta=text] {#1};}

% Short panel titles; the captions carry the full read-set descriptions.
\newcommand{\figTitle}[1]{\csname figtitle:#1\endcsname}
\expandafter\def\csname figtitle:human_short\endcsname{Short reads}
\expandafter\def\csname figtitle:human_pacbio\endcsname{PacBio HiFi}
\expandafter\def\csname figtitle:human_ont_modern\endcsname{Nanopore}
\expandafter\def\csname figtitle:guccione_skin\endcsname{Skin metagenome}
\expandafter\def\csname figtitle:simshort\endcsname{Sim.\ short reads}
\expandafter\def\csname figtitle:simpacbio\endcsname{Sim.\ PacBio HiFi}
\expandafter\def\csname figtitle:simont\endcsname{Sim.\ Nanopore}
\expandafter\def\csname figtitle:cap1\endcsname{\textit{N} = 1}
\expandafter\def\csname figtitle:cap10\endcsname{\textit{N} = 10}
\expandafter\def\csname figtitle:cap100\endcsname{\textit{N} = 100}

% \figPanelBody{fig}{panel}{link layouts}{keys}: everything drawn inside one
% panel, after its \nextgroupplot.
\newcommand{\figPanelBody}[4]{%
  \figPlotLinks{data/#1/#2.dat}{#3}%
  \figPlotKeys{data/#1/#2.dat}{#4}%
  \figPlotLabels{data/#1/#2-labels.dat}%
  \figPanelTitle{\figTitle{#2}}}

% Legend groups. An empty-legend entry names the tool before its modes.
\newcommand{\figLegendGroupName}[1]{\addlegendimage{empty legend}\addlegendentry{\textit{#1}}}
\newcommand{\figLegendKey}[2]{\addlegendimage{ms key/#1}\addlegendentry{#2}}
\newcommand{\figLegendMS}{%
  \figLegendGroupName{ms2:}%
  \figLegendKey{lcp}{min}\figLegendKey{psi}{psi}\figLegendKey{dual}{dual}%
  \figLegendGroupName{ropebwt3:}%
  \figLegendKey{rl1}{-l1}\figLegendKey{rl19}{-l19*}\figLegendKey{rl25}{-l25*}%
  \figLegendKey{rl30}{-l30*}\figLegendKey{rl40}{-l40*}}
\newcommand{\figLegendOracle}{%
  \addlegendimage{fig hollow={figInkTwo}{square}, mark size=1.7pt}\addlegendentry{oracle}}
\newcommand{\figLegendFtab}{%
  \addlegendimage{fig filled={figInkTwo}{*}}\addlegendentry{ordinary}%
  \addlegendimage{fig filled={figInkTwo}{diamond*}, mark size=2.3pt}\addlegendentry{ftab}}
\newcommand{\figLegendOthers}{%
  \figLegendKey{moni}{MONI}\figLegendKey{phoni}{PHONI}\figLegendKey{augphoni}{AugPHONI}}
\newcommand{\figLegendLocate}{%
  \figLegendGroupName{ms2:}%
  \figLegendKey{psi}{psi}\figLegendKey{dual}{dual}%
  \figLegendGroupName{ropebwt3 -l1:}%
  \figLegendKey{s11}{-s11}\figLegendKey{s8}{-s8}\figLegendKey{s6}{-s6}}

% Shared x label and legend placed relative to a row of panels.
% \figRowFurniture{first}{last}{x label}{legend name}
\newcommand{\figRowFurniture}[4]{%
  \node[font=\figLabelFont, text=figInkTwo, anchor=north]
    at ($(#1.south west)!0.5!(#2.south east)+(0,-8pt)$) {#3};
  \node[anchor=south, inner sep=0pt]
    at ($(#1.north west)!0.5!(#2.north east)+(0,2pt)$) {\pgfplotslegendfromname{#4}};}

\newcommand{\figMemLabel}{Peak resident memory (GiB)}
\newcommand{\figSpeedLabel}{Speed (Mbases/s)}
\newcommand{\figLocLabel}{Locations/s}

% Fig Ops: bar panels for time and operations per base (data/fig-ops/).
% Bars are drawn per layout key with zero bar shift so each sits on its own x.
% An ops bar is the full total in the layout color with the LF segment (one
% step per base, which is also what the oracle performs) painted over it as a
% 45 percent tint; the white rim leaves a surface gap between the segments.
\newlength{\figOpsPanelW}\setlength{\figOpsPanelW}{1.25in}
\newlength{\figOpsCountPanelW}\setlength{\figOpsCountPanelW}{1.0in}
\pgfplotsset{
  ops axis/.style={
    fig axis, width=\figOpsPanelW, xmin=-0.6, xmax=3.6, ymin=0,
    xtick={0,1,2,3}, xticklabels={min,psi,phi,dual}, major x tick style={draw=none},
    xticklabel style={font=\figLabelFont, text=figInkTwo},
    % Numeric tick labels outside math mode, so they take the sans tick font.
    yticklabel style={/pgf/number format/assume math mode=true},
  },
  % Operation panels leave out ms2-min (keyed lcp), whose only added work is a short LCP scan.
  ops count axis/.style={ops axis, width=\figOpsCountPanelW, xmin=0.4, xmax=3.6,
    xtick={1,2,3}, xticklabels={psi,phi,dual}},
  ops bar/.style={ybar, bar shift=0pt, bar width=0.62, forget plot, draw=white, line width=0.5pt, fill=#1},
  ops tick/.style={only marks, forget plot, mark=-, mark size=8.6pt,
    mark options={draw=figInk, line width=1.0pt, line cap=butt}},
  ops labels/.style={only marks, mark=none, forget plot, point meta=explicit symbolic,
    nodes near coords, nodes near coords style={font=\figPointFont, text=figInkTwo, inner sep=0.6pt,
      anchor=south, yshift=0.6pt}},
  ops legend bar/.style={legend image code/.code={\fill[#1] (0pt,-2.6pt) rectangle (11pt,2.6pt);}},
}
\newcommand{\figOpsColor}[1]{\csname figopscolor:#1\endcsname}
\expandafter\def\csname figopscolor:lcp\endcsname{msLcp}
\expandafter\def\csname figopscolor:psi\endcsname{msPsi}
\expandafter\def\csname figopscolor:dual\endcsname{msDual}
\expandafter\def\csname figopscolor:phi\endcsname{msPhi}
% \figOpsTime{file}: ordinary time bars, oracle ticks and bar labels.
\newcommand{\figOpsTime}[1]{%
  \pgfplotsinvokeforeach{lcp,psi,dual,phi}{%
    \addplot[ops bar=\figOpsColor{##1}, discard if not={key}{##1}] table[x=x, y=ns] {#1};}%
  \addplot[ops tick] table[x=x, y=oracle] {#1};
  \addplot[ops labels] table[x=x, y=ns, meta=nslab] {#1};}
% \figOpsOps{file}: stacked operation bars and total labels.
\newcommand{\figOpsOps}[1]{%
  \pgfplotsinvokeforeach{psi,dual,phi}{%
    \addplot[ops bar=\figOpsColor{##1}, discard if not={key}{##1}] table[x=x, y=total] {#1};
    \addplot[ops bar=\figOpsColor{##1}!45, discard if not={key}{##1}] table[x=x, y=lf] {#1};}%
  \addplot[ops labels, discard if={key}{lcp}] table[x=x, y=total, meta=totlab] {#1};}
\newcommand{\figLegendOps}{%
  \addlegendimage{ops legend bar=msLcp}\addlegendentry{min}%
  \addlegendimage{ops legend bar=msPsi}\addlegendentry{psi}%
  \addlegendimage{ops legend bar=msPhi}\addlegendentry{phi}%
  \addlegendimage{ops legend bar=msDual}\addlegendentry{dual}%
  \addlegendimage{line legend, draw=figInk, line width=1.0pt}\addlegendentry{oracle time}%
  \addlegendimage{ops legend bar=figInkTwo!45}\addlegendentry{LF steps}%
  \addlegendimage{ops legend bar=figInkTwo}\addlegendentry{Recovery ($\psi$ \& $\phi$) steps}}
% Bar panels put the title above the plot, clear of the bar labels.
\newcommand{\figOpsTitle}[1]{%
  \node[anchor=south west, inner sep=1.5pt, font=\figTitleFont, text=figInk]
    at (rel axis cs:0,1) {#1};}
\newcommand{\figOpsTimeLabel}{ns per base (query time)}
\newcommand{\figOpsLabel}{Operations per base}
\expandafter\def\csname figtitle:mtb152\endcsname{mtb152}
\expandafter\def\csname figtitle:human472\endcsname{HPRC2}
